\chapter{Introduction}
\label{ch:intro}

\section{Motivation and Problem Statement}

Deciding what to cook from whatever ingredients happen to be available is a
common practical problem. Users typically know what is in their kitchen, but
translating that into a concrete, safe, and personalized meal requires browsing
multiple recipe websites, manually filtering results, and adapting instructions
to fit dietary restrictions or allergies. This process is repetitive and
effortful, yet users still often lack a simple tool that accepts a list of
ingredients and returns a structured recipe adapted to personal constraints.

A related inconvenience arises with packaged food products. Nutritional
information is printed on product labels in a format that is easy to read in
the moment but difficult to store, compare, or reuse digitally. Users who
want to track or inspect the nutritional content of packaged goods have no
practical way to capture that data without manually transcribing it, which
is both time-consuming and error-prone.

This thesis addresses both problems through a single integrated application.
The core feature allows users to submit a custom ingredient list along with
optional dietary preferences and allergy constraints, and receive a structured
recipe generated specifically from that input. A secondary feature accepts
an image of a food product label, extracts the printed nutritional text
automatically, and allows users to store the result in a structured digital form. Together,
these features form a coherent tool for ingredient-aware cooking decisions
and nutritional data capture within a single personalized cooking
assistant. 

\section{Goals and Objectives}

The main goal of the thesis is to deliver a working web application that
generates personalized recipes from user-provided ingredients and supports
OCR-based nutrition label processing. The project commits to four concrete
and verifiable objectives.

The first objective is to implement an ingredient-based recipe generation
workflow. The system accepts a list of ingredients and optional dietary
preferences and allergies, validates the request, and returns a structured
recipe containing a title, ingredient list, preparation steps, nutritional
estimate, and system warnings. To improve the relevance of generated
results, the workflow uses retrieval-augmented generation: before producing
a recipe, the system retrieves relevant entries from an indexed recipe corpus
and uses them as contextual grounding for the language model.

The second objective is to implement user-specific personalization and
persistence. Authenticated users can create accounts, save dietary
preferences and allergies, store generated recipes, and access both saved
recipes and their full generation history. This ensures the application
supports repeated and personalized use rather than one-time generation.

The third objective is to implement OCR-based nutrition label extraction.
The system accepts uploaded label images, validates and preprocesses them,
extracts text through an OCR provider, transforms the raw output into
structured nutritional fields, and allows authenticated users to save their
labels and extracted data.

The fourth objective is to implement and document the backend and data model
that support these workflows. This covers the system architecture, request
handling, authentication, persistent storage, media handling, and the
service boundaries required to coordinate recipe generation, OCR
processing, and history management.

\section{Scope and Constraints}

The scope of the thesis covers the implemented \emph{AI Cookbook with
Ingredient-Based Search} application and all features available in the final
system: ingredient-based recipe generation, preference- and allergy-aware
personalization, OCR-based nutrition label extraction, user authentication,
and storage of recipe and label history. It also covers the technical
documentation of the system, including the application architecture, main
request flows, persistent data model, and testing material.

The thesis does not aim to provide a general introduction to artificial
intelligence, nutrition science, or OCR theory beyond what is necessary to
explain the implemented solution.

Several constraints apply. Recipe generation, nutritional structuring, and
image generation depend on external AI services, meaning that response
quality and availability are not fully controlled by the local application.
OCR accuracy depends on the quality and readability of uploaded images.
Dietary preferences, allergy-related handling, and nutritional values
displayed by the system, whether extracted from labels or estimated for
generated recipes, are intended as application output for user guidance and
do not constitute certified dietary or medical advice.

The project also does not aim to provide production-grade infrastructure,
clinical nutritional analysis, or guaranteed factual correctness for every
AI-generated or OCR-derived value. Its purpose is to demonstrate a working
and well-documented implementation of the designed workflows within the
scope of a bachelor thesis.

The project is also constrained by its scope as a bachelor thesis. The
system was designed to be functionally complete and documentable within
limited development time. The thesis therefore focuses on the implemented
features and their software engineering realization rather than
production-scale concerns such as distributed deployment or expert-level
nutritional analytics. Additional feature expansion and broader
implementation improvements are therefore treated as future work.

\section{Thesis Roadmap}

Chapter 2 presents the user documentation and describes how the application
is configured, launched, and used from the perspective of its intended users.

Chapter 3 presents the developer documentation. It is divided into the
design, implementation, and testing sections, covering the software
requirements, system architecture, data model, implementation details,
and validation results of the application.

Chapter 4 concludes the thesis by summarizing the achieved results,
discussing the limitations of the current implementation, and outlining
directions for future improvement.
